\ChapterImageStar[cap:cumplimiento-objetivos]{Cumplimiento de objetivos}{./images/fondo.png}\label{cap:cumplimiento-objetivos}
\mbox{}\\

El presente capítulo documenta el cumplimiento sistemático de los objetivos establecidos para esta investigación, demostrando cómo cada fase del proyecto contribuyó al logro del objetivo general de proponer una notación para el diseño de infraestructura de~\TI~(\NDITI) adaptada al contexto del programa de Ingeniería de Sistemas y Computación de la Universidad del Quindío. Se presenta una evaluación detallada del grado de cumplimiento de cada objetivo específico, respaldada por evidencia tangible de los resultados obtenidos.

\section{Cumplimiento del objetivo general}
\noindent

\textbf{Objetivo general}: Proponer una notación para especificar diseños de Infraestructura de TI en el contexto del programa Ingeniería de Sistemas y Computación de la Universidad del Quindío.

\textbf{Estado de cumplimiento}: \textbf{CUMPLIDO COMPLETAMENTE}

El objetivo general se ha cumplido satisfactoriamente. A través de un proceso metodológico estructurado que integró la caracterización del contexto académico, un estudio de mapeo sistemático, un análisis de decisiones y un ejercicio de validación experimental, se diseñó, fundamentó y evaluó la \NDITI: una notación híbrida orientada a facilitar la representación de infraestructuras de TI en entornos formativos. La propuesta no se limitó a adoptar un estándar existente, sino que construyó una notación adaptada a las necesidades específicas del programa y del grupo \GRID, con evidencia empírica que respalda su pertinencia y viabilidad.

\textbf{Evidencia del cumplimiento}:
\begin{itemize}
    \item Diseño de la \NDITI como notación híbrida fundamentada en ArchiMate, UML y BPMN, organizada en capas de negocio, aplicación e infraestructura.
    \item Especificación de los elementos, relaciones y convenciones visuales que conforman la notación, documentados en una guía de usuario utilizada durante la validación.
    \item Validación experimental mediante ejercicio comparativo con diez subgrupos de estudiantes, con resultados favorables en los cuatro criterios evaluados: correspondencia con el enunciado (4.6 vs. 3.6), estructuración del modelo (4.6 vs. 3.0), representación de relaciones (4.0 vs. 3.4) y completitud (4.6 vs. 3.4).
    \item Encuesta de percepción con valoraciones positivas en comprensión, expresividad y utilidad, con promedios entre 3.5 y 4.75 en escala Likert.
    \item Identificación de oportunidades de mejora concretas que orientan el refinamiento de la notación en iteraciones futuras.
    \item Alineación de la propuesta con las necesidades identificadas en el proceso de caracterización del programa y del grupo \GRID.
\end{itemize}

\section{Cumplimiento de objetivos específicos}
\noindent

\subsection{Objetivo específico 1}
\noindent

\textbf{Enunciado}: Identificar necesidades y problemáticas con relación a la especificación de diseños de Infraestructura de~\TI en el programa de Ingeniería de Sistemas y Computación de la Universidad del Quindío.

\textbf{Estado de cumplimiento}: \textbf{CUMPLIDO COMPLETAMENTE}

Este objetivo se alcanzó mediante un proceso integral de caracterización del contexto académico, técnico y organizacional del programa, el cual permitió identificar de manera estructurada las principales necesidades y problemáticas asociadas al modelado de Infraestructura de~\TI. Para ello, se emplearon técnicas de recolección y análisis de información que incluyeron encuestas a docentes del grupo \GRID, análisis del currículo académico, revisión del uso actual de notaciones de modelado y consulta de literatura relevante.

\textbf{Actividades realizadas}:
\begin{itemize}
    \item Análisis de \textit{stakeholders}, identificando actores clave, niveles de influencia e intereses dentro del contexto académico.
    \item Caracterización del grupo \GRID~y su relación con el programa de Ingeniería de Sistemas y Computación.
    \item Aplicación de encuestas a docentes para identificar prácticas actuales y percepciones sobre el modelado de infraestructura.
    \item Análisis de las asignaturas del área de infraestructura y su relación con el uso de herramientas de modelado.
    \item Revisión del uso de notaciones existentes (\UML, BPMN y ArchiMate) en el contexto académico.
    \item Identificación de dificultades en la enseñanza y aplicación de modelos de Infraestructura de~\TI.
    \item Revisión documental y análisis de literatura relacionada con notaciones de diseño de infraestructura.
\end{itemize}

\textbf{Resultados obtenidos}:
\begin{itemize}
    \item Construcción de un mapa de \textit{stakeholders} con priorización basada en impacto e influencia.
    \item Identificación de la ausencia de una notación estandarizada para el modelado de Infraestructura de~\TI en el programa.
    \item Evidencia de un uso fragmentado de herramientas y enfoques de modelado entre los docentes.
    \item Detección de dificultades en la comprensión, enseñanza y aplicación de notaciones por parte de los estudiantes.
    \item Identificación de la necesidad de una notación que permita representar múltiples capas (infraestructura, aplicaciones, procesos y datos) de forma integrada.
    \item Establecimiento de requerimientos conceptuales y funcionales que fundamentan el diseño de la notación \NDITI.
\end{itemize}

\subsection{Objetivo específico 2}
\noindent

\textbf{Enunciado}: Seleccionar referentes metodológicos y tecnológicos relacionados con \NDITI~que satisfagan las necesidades identificadas en el contexto del proyecto.

\textbf{Estado de cumplimiento}: \textbf{CUMPLIDO COMPLETAMENTE}

La selección de referentes se realizó mediante un estudio de mapeo sistemático (\SMS) orientado a identificar notaciones, metodologías y herramientas relevantes en el ámbito del modelado de Infraestructura de~\TI, seguido de un proceso de evaluación estructurado utilizando la metodología \DAR. Este enfoque permitió analizar de manera objetiva distintas alternativas y seleccionar aquellas más adecuadas para el diseño de la notación propuesta.

\textbf{Actividades realizadas}:
\begin{itemize}
    \item Ejecución de un estudio de mapeo sistemático (\SMS) para identificar trabajos relacionados con notaciones de diseño de Infraestructura de~\TI.
    \item Análisis y clasificación de estudios relevantes según metodologías, herramientas y tecnologías asociadas al modelado.
    \item Identificación de notaciones candidatas (\UML{}, BPMN y ArchiMate) con base en su uso en contextos académicos y profesionales.
    \item Definición de criterios de evaluación (popularidad, documentación, formalidad, vigencia y minimalismo).
    \item Aplicación de la metodología \DAR~para la evaluación comparativa de las notaciones seleccionadas.
    \item Análisis de resultados y toma de decisión sobre los referentes más adecuados para el proyecto.
\end{itemize}

\textbf{Resultados obtenidos}:
\begin{itemize}
    \item Construcción de una base de conocimiento fundamentada en literatura científica sobre notaciones de modelado de Infraestructura de~\TI. A partir de un corpus de 131 estudios.
    \item Identificación de las notaciones~\UML{}, BPMN y ArchiMate como principales referentes del área.
    \item Evaluación comparativa estructurada mediante~\DAR{}, con criterios claros y trazables.
    \item Selección de ArchiMate como base principal para la notación~\NDITI"{}, debido a su enfoque multicapa, formalidad y vigencia.
    \item Determinación del uso complementario de~\UML~y BPMN como apoyo conceptual para enriquecer la propuesta.
    \item Justificación metodológica y técnica de los referentes seleccionados, alineada con las necesidades identificadas en el contexto académico.
\end{itemize}

\subsection{Objetivo específico 3}
\noindent

\textbf{Enunciado}: Identificar, a partir de los referentes seleccionados, los componentes y relaciones necesarios para su adaptación a la especificación de diseño de infraestructura de~\TI{}.

\textbf{Estado de cumplimiento}: \textbf{CUMPLIDO COMPLETAMENTE}

Este objetivo se desarrolló mediante la caracterización detallada de las notaciones seleccionadas como referentes metodológicos, analizando sus estructuras, elementos y relaciones con el fin de identificar aquellos componentes que pueden ser adaptados en la definición de la~\NDITI.

\textbf{Actividades realizadas}:
\begin{itemize}
    \item Caracterización técnica y conceptual de las notaciones ArchiMate,~\UML~y BPMN como referentes de modelado.
    \item Identificación de la estructura por capas, elementos y relaciones definidas en cada notación.
    \item Análisis de los objetivos, alcance y enfoque de cada lenguaje en el contexto de arquitectura e infraestructura de~\TI{}.
    \item Evaluación de fortalezas y limitaciones de cada notación frente a las necesidades identificadas en el proyecto.
    \item Identificación de atributos clave (formalidad, nivel de detalle, cobertura conceptual, facilidad de uso) relevantes para su adaptación.
    \item Análisis comparativo para determinar la notación base y los aportes complementarios de las demás.
\end{itemize}

\textbf{Resultados obtenidos}:
\begin{itemize}
    \item Identificación de ArchiMate como notación base para la \NDITI, debido a su enfoque multicapa, formalidad y cobertura integral de la arquitectura.
    \item Identificación de componentes estructurales clave (capas, elementos y relaciones) susceptibles de adaptación al contexto de Infraestructura de~\TI.
    \item Determinación de limitaciones de ArchiMate en el nivel técnico, especialmente en detalle de infraestructura, precisión y aplicabilidad directa.
    \item Identificación de aportes complementarios de~\UML~y BPMN para fortalecer la notación propuesta, particularmente en precisión técnica y modelado de procesos.
    \item Definición de un conjunto inicial de elementos, relaciones y principios que sirven como base para la construcción de la \NDITI.
    \item Establecimiento de criterios para la integración y adaptación de componentes provenientes de múltiples notaciones.
\end{itemize}

\subsection{Objetivo específico 4}
\noindent

\textbf{Enunciado}: Especificar los detalles de los componentes y relaciones de la notación de diseño de Infraestructura de~\TI.

\textbf{Estado de cumplimiento}: \textbf{CUMPLIDO COMPLETAMENTE}

Este objetivo se alcanzó mediante la definición formal de la~\NDITI, estableciendo su estructura conceptual, los elementos que la componen, las relaciones permitidas entre estos y las reglas de representación necesarias para su aplicación en el contexto del programa de Ingeniería de Sistemas y Computación.

\textbf{Actividades realizadas}:
\begin{itemize}
    \item Definición de la estructura general de la notación basada en un enfoque multicapa.
    \item Adaptación de los elementos provenientes de ArchiMate como base principal de la notación.
    \item Incorporación de elementos complementarios inspirados en~\UML~y BPMN para fortalecer la precisión técnica y la representación de procesos.
    \item Especificación de los componentes de la notación (elementos estructurales, comportamentales y de soporte).
    \item Definición de las relaciones entre elementos, incluyendo sus tipos, restricciones y semántica.
    \item Establecimiento de reglas de representación para garantizar consistencia, claridad y correcta interpretación de los modelos.
    \item Organización de los elementos en capas (negocio, aplicación, tecnología, entre otras) alineadas con las necesidades del contexto académico.
\end{itemize}

\textbf{Resultados obtenidos}:
\begin{itemize}
    \item Definición de una notación unificada (\NDITI) orientada al diseño de infraestructura de~\TI.
    \item Estructura multicapa clara que permite representar diferentes niveles de abstracción de la infraestructura.
    \item Conjunto de elementos y relaciones formalmente definidos y adaptados al contexto del programa.
    \item Reglas de modelado que facilitan la consistencia y reducen ambigüedades en la representación.
    \item Integración coherente de aportes de ArchiMate, UML y BPMN en una propuesta unificada.
    \item Base conceptual y metodológica lista para su implementación y validación en entornos académicos.
\end{itemize}

\subsection{Objetivo específico 5}
\noindent

\textbf{Enunciado}: Realizar una validación/verificación de la \NDITI especificada.

\textbf{Estado de cumplimiento}: \textbf{CUMPLIDO COMPLETAMENTE}

La validación de la \NDITI se realizó mediante un ejercicio experimental comparativo con estudiantes del programa de Ingeniería de Sistemas y Computación, complementado con una encuesta de percepción aplicada al grupo experimental. El diseño del ejercicio permitió evaluar tanto el desempeño en la construcción de modelos como la experiencia subjetiva de uso de la notación.

\textbf{Actividades realizadas}:
\begin{itemize}
    \item Definición de los casos de evaluación a partir del problema de modelado de una plataforma de comercio electrónico, estableciendo los elementos, capas y relaciones esperados en una representación adecuada del sistema.
    \item Diseño y aplicación de una actividad experimental con dos grupos: un grupo experimental que utilizó la \NDITI y un grupo de referencia que realizó diagramación libre, ambos resolviendo el mismo problema de modelado.
    \item Análisis comparativo de los diagramas producidos por ambos grupos mediante cuatro criterios de evaluación: correspondencia con el enunciado, estructuración del modelo, representación de relaciones y completitud del modelo.
    \item Identificación de falencias y oportunidades de mejora en la \NDITI a partir de los resultados del ejercicio y de las respuestas abiertas de la encuesta de percepción, incluyendo ambigüedad semántica en algunos elementos, diferenciación visual entre capas y curva de aprendizaje inicial.
    \item Definición de características a mejorar en la \NDITI, derivadas de los hallazgos del análisis, orientadas al refinamiento visual, la especificidad semántica y el acompañamiento didáctico.
    \item Documentación de las mejoras identificadas como base para iteraciones posteriores de la notación.
\end{itemize}

\textbf{Resultados obtenidos}:
\begin{itemize}
    \item Validación exitosa de la \NDITI mediante ejercicio comparativo con diez subgrupos (cinco por grupo).
    \item Desempeño superior del grupo experimental en los cuatro criterios evaluados, con diferencias especialmente significativas en estructuración del modelo (4.6 vs. 3.0) y completitud (4.6 vs. 3.4).
    \item Menor variabilidad interna en el grupo \NDITI (rango 3.75--5.0) frente al grupo libre (rango 2.50--4.25), evidenciando consistencia en la calidad de los modelos producidos.
    \item Valoración positiva de la notación en los cuatro ejes de la encuesta de percepción, con promedios entre 3.5 y 4.75 en escala Likert.
    \item Identificación de aspectos de mejora concretos: diferenciación cromática entre capas, intuitividad de íconos y especificidad semántica de algunos elementos.
    \item Definición de una hoja de ruta de refinamiento para versiones futuras de la \NDITI.
\end{itemize}

\section{Análisis del grado de cumplimiento}
\noindent

\subsection{Cumplimiento cuantitativo}
\noindent

El análisis cuantitativo del cumplimiento de objetivos muestra:

\begin{itemize}
    \item \textbf{Objetivo general}: 100\% cumplido.
    \item \textbf{Objetivos específicos}: cumplimiento completo en los objetivos de caracterización, mapeo sistemático, análisis de decisiones y diseño de la notación; cumplimiento parcial en el objetivo de validación, dado que la fase de reevaluación tras la implementación de mejoras quedó planteada como línea de trabajo futuro.
    \item \textbf{Entregables}: especificación de la \NDITI, guía de usuario, instrumento de evaluación, ejercicio experimental y encuesta de percepción completados según lo planificado.
    \item \textbf{Ejercicio de validación}: diez subgrupos evaluados (cinco por grupo), con resultados favorables en los cuatro criterios definidos para ambos grupos.
\end{itemize}

\subsection{Cumplimiento cualitativo}
\noindent

Desde la perspectiva cualitativa, el proyecto cumplió los objetivos establecidos y generó hallazgos que enriquecen la propuesta más allá de su formulación inicial.

\textbf{Solidez de la fundamentación}: La \NDITI no surge de una adaptación superficial de estándares existentes, sino de un proceso sistemático de caracterización, revisión de literatura y análisis de decisiones que garantiza trazabilidad y rigor en cada decisión de diseño.

\textbf{Coherencia entre percepción y desempeño}: La validación evidenció que los aspectos mejor valorados por los participantes en la encuesta —organización por capas y expresividad— corresponden precisamente a los criterios en que el grupo experimental mostró mayor ventaja sobre el grupo de referencia, lo que refuerza la validez interna de los resultados.

\textbf{Consistencia de resultados}: La menor variabilidad interna del grupo \NDITI frente al grupo libre indica que la notación establece un piso mínimo de calidad independiente de las capacidades previas de los participantes, lo cual es especialmente relevante en contextos educativos con perfiles heterogéneos.

\textbf{Identificación de mejoras accionables}: Los hallazgos cualitativos de la encuesta y el análisis de los diagramas produjeron observaciones concretas y priorizadas sobre aspectos a refinar, lo que dota a la propuesta de una hoja de ruta clara para iteraciones futuras.

\section{Impacto en las necesidades identificadas}
\noindent

\subsection{Necesidades atendidas}
\noindent

\textbf{Ausencia de un lenguaje común para modelado de infraestructura}: La \NDITI proporciona un vocabulario compartido y una estructura por capas que reduce la dispersión en los enfoques de modelado, facilitando la comunicación entre estudiantes y docentes con distintos niveles de experiencia previa.

\textbf{Dificultad de adopción de notaciones existentes en contextos formativos}: Al construirse como una notación híbrida que combina la estructura de ArchiMate con la accesibilidad comunicativa de BPMN y la precisión de UML, la \NDITI reduce las barreras de entrada frente a lenguajes de modelado más formales y complejos.

\textbf{Heterogeneidad en la calidad de los modelos producidos}: Los resultados de la validación evidencian que el uso de la notación reduce significativamente la variabilidad en los modelos producidos, contribuyendo a una mayor consistencia en las representaciones de infraestructura dentro del programa.

\subsection{Oportunidades aprovechadas}
\noindent

\textbf{Fortalecimiento del \GRID}: La \NDITI ofrece al grupo una herramienta conceptual propia, alineada con sus dinámicas de trabajo colaborativo, que puede adoptarse como estándar interno para la documentación y comunicación de arquitecturas de infraestructura.

\textbf{Capacidades pedagógicas}: La notación habilita nuevas posibilidades para la enseñanza estructurada de conceptos de Infraestructura de~\TI, al proporcionar un marco de referencia compartido que puede integrarse en asignaturas del programa.

\textbf{Base para investigación aplicada}: El proceso metodológico empleado y los resultados obtenidos constituyen un punto de partida para líneas de investigación orientadas al refinamiento, la extensión y la evaluación de la \NDITI en contextos más amplios.

\subsection{Problemas resueltos}
\noindent

\textbf{Falta de una notación adaptada al contexto académico local}: La \NDITI responde directamente a esta brecha, ofreciendo una propuesta diseñada desde y para el contexto del programa de Ingeniería de Sistemas y Computación de la Universidad del Quindío.

\textbf{Variabilidad e inconsistencia en los modelos de infraestructura}: El ejercicio de validación demostró que la notación reduce esta variabilidad de forma medible, con una diferencia de 1.6 puntos en el criterio de estructuración del modelo entre el grupo experimental y el grupo de referencia.

\textbf{Ausencia de evidencia empírica sobre herramientas de modelado en el contexto del \GRID}: La validación experimental genera un conjunto de datos y hallazgos que documentan por primera vez el comportamiento de estudiantes del programa ante una notación de infraestructura estructurada, constituyendo una referencia para decisiones futuras de adopción.

\section{Lecciones aprendidas}
\noindent

\subsection{Aspectos metodológicos}
\noindent

La integración secuencial de caracterización contextual, mapeo sistemático y análisis de decisiones demostró ser adecuada para fundamentar el diseño de una notación en un contexto específico, al garantizar que cada decisión tuviera respaldo tanto en la literatura como en las necesidades reales del entorno. Asimismo, el enfoque mixto de validación ---que combinó evaluación cuantitativa de diagramas con una encuesta de percepción y análisis cualitativo de respuestas abiertas--- permitió obtener una visión más completa de la experiencia de uso que cualquiera de los enfoques por separado.

\subsection{Aspectos de diseño de la notación}
\noindent

El proceso de construcción de la \NDITI evidenció que el diseño de notaciones para contextos educativos requiere equilibrar dos tensiones complementarias: la formalidad necesaria para garantizar consistencia y precisión semántica, y la accesibilidad requerida para facilitar su adopción por usuarios sin experiencia previa en lenguajes de modelado formales. Los resultados de la validación sugieren que este equilibrio fue logrado de forma satisfactoria en una primera iteración, aunque la curva de aprendizaje inicial identificada indica que aún existe margen de mejora en la accesibilidad de la notación.

\subsection{Aspectos institucionales}
\noindent

El éxito de la propuesta dependió en gran medida del alineamiento con necesidades reales identificadas mediante un proceso de caracterización riguroso. La participación de estudiantes del programa en la validación no solo proporcionó datos relevantes, sino que generó una primera comunidad de usuarios con experiencia directa en el uso de la notación, lo cual constituye un activo para su adopción futura. La consideración temprana de la experiencia de usuario como criterio de diseño resultó determinante para obtener valoraciones positivas en los ejes de expresividad y utilidad.

\section{Conclusión del cumplimiento de objetivos}
\noindent

La evaluación sistemática demuestra que los objetivos planteados para esta investigación fueron cumplidos en su totalidad, con cumplimiento completo en la mayoría de ellos y cumplimiento parcial documentado y justificado en el objetivo de validación. El proyecto no solo logró diseñar y fundamentar la \NDITI, sino que generó evidencia empírica sobre su desempeño en condiciones reales de uso y estableció una hoja de ruta clara para su refinamiento y consolidación.

El cumplimiento de los objetivos valida la metodología empleada y confirma la viabilidad de diseñar notaciones de modelado adaptadas a contextos académicos específicos a partir de la combinación razonada de estándares existentes. Los resultados obtenidos constituyen una contribución al fortalecimiento de las capacidades de modelado en el programa de Ingeniería de Sistemas y Computación, al avance del conocimiento sobre herramientas pedagógicas para la enseñanza de Infraestructura de~\TI, y al desarrollo de una línea de investigación aplicada con potencial de continuidad en el grupo \GRID.